\chapter{Expected Outcome}

The expected outcome of the Farming Intelligence AI project is a user-friendly agricultural decision-support system that assists farmers in making informed decisions using soil, weather, crop, and market information. The system integrates IoT-based soil monitoring, machine learning, agricultural advisory, explainable AI, market intelligence, and a retrieval-augmented generation (RAG) assistant into a single platform. The crop recommendation module is expected to identify suitable crops based on parameters such as nitrogen, phosphorus, potassium, temperature, humidity, pH, and rainfall, while live soil sensor data provides updated field conditions for more relevant recommendations. The fertilizer and irrigation advisory module is expected to provide practical guidance based on soil nutrients, moisture, pH, and weather conditions, helping farmers manage fertilizers and water more efficiently. The Explainable AI module is expected to improve transparency by using SHAP and LIME to identify the important factors influencing crop predictions. The market intelligence module is expected to provide market trends, estimated future prices, and profitability information to support better crop selection. The RAG-based agricultural assistant is expected to provide context-based agricultural information from the prepared knowledge base, while multilingual support improves accessibility for farmers. Overall, the project is expected to provide an integrated smart farming platform that supports data-driven crop selection, resource management, agricultural guidance, and market-based decision-making.

The system is also expected to provide a simple and accessible interface through which farmers can view their farm information, live soil conditions, recommendations, advisory results, market information, and AI-generated agricultural guidance. By bringing these functionalities together in a single platform, the project aims to reduce the complexity of agricultural decision-making and support farmers in adopting more informed and efficient farming practices.